\section{Формулировки определений и утверждений}

Логические символы:
\begin{description}
    \item[] $\forall$ --- для любого (квантор всеобщности)
    \item[] $\exists$--- существует (квантор существования)
    \item[] $\colon$ --- такой, что
    \item[] $\to$ --- выполняется
    \item[] $A \implies B$ --- из условия $A$ следует утверждение $B$ ($B$ является следствием $A$)
    \item[] $A \iff B$ --- условия $A$ и $B$ равносильны
    \item[] $\neg A$ --- отрицание условия $A$
    \item[] $A \land B$ --- логическое пересечение (одновременно выполнены оба условия)
    \item[] $A \lor B$ --- логическое объединение (выполнено хотя бы одно из условий)
\end{description}

$A \implies B$:
\begin{enumerate}
    \item $A$ является \textit{достаточным} условием для $B$.
    \item $B$ является \textit{необходимым} условием для $A$.
\end{enumerate}

$A \iff B$: условие $A$ является \textit{критерием} выполнения условия $B$.

Символ $\to$ используется внутри определений и утверждений.

Символ $\implies$ используется между утверждениями.

\begin{definition}{Отрицание утверждения}{}
    Правило построения отрицания утверждения:
    \begin{itemize}
        \item Квантор $\forall$ заменяется на $\exists$, и наоборот
        \item Условие $A$ в конце утверждения заменяется на его отрицание $\neg A$
        \item После квантора $\exists$ ставится двоеточие $\colon$
        \item После квантора $\forall$ ставится стрелка $\to$
    \end{itemize}
\end{definition}

\begin{example}{}{}
    Пусть $A(x)$ --- некоторое условие на элементы $x \in X$. Тогда:
    \begin{enumerate}
        \item $\neg (\forall x \in X \to A(x)) \iff \exists x \in X \colon \neg A(x)$
        \item $\neg (\exists x \in X \colon A(x)) \iff \forall x \in X \to \neg A(x)$
    \end{enumerate}
\end{example}

Применение знака $=$:
\begin{enumerate}
    \item $B \coloneqq F(A)$ --- определение объекта $B$ с помощью объекта $A$.
    \item $f(x) \equiv g(x)$ при $x \in X$ --- тождественное равенство отображений на множестве $X$,
        где $X \subset \D(f)$ и $X \subset \D(g)$.
\end{enumerate}

\begin{example}{}{}
    \begin{enumerate}
        \item По определению функция $\tg(x) \coloneqq \sin(x) / \cos(x)$
        \item $|x| \equiv x$ при $x \ge 0$
    \end{enumerate}
\end{example}

\begin{theorem}{Метод математической индукции}{}
    Пусть в условие утверждения входит параметр $n \in \N$
    и требуется доказать его справедливость для всех $n \in \N$.
    Для этого достаточно установить, что:
    \begin{enumerate}
        \item Утверждение верно при $n = 1$ (\textit{база индукции})
        \item Если $\forall n \in \N$ утверждение верно,
            то оно верно и для $n+1$ (\textit{индукционный переход})
    \end{enumerate}
\end{theorem}

\begin{theorem}{Метод доказательства от противного}{}
    Для доказательства утверждение $A$ можно поступить следующим образом:
    \begin{enumerate}
        \item Формулируется отрицание $\neg A$ доказываемого утверждения $A$
            (\textit{предположение, гипотеза})
        \item Из предположения $\neg A$ логически выводится
            заведомо ложное утверждение $B$ (\textit{противоречие})
    \end{enumerate}
    Полученное противоречие показывает, что предположение $\neg A$ было неверным,
    и поэтому верно его отрицание $\neg \neg A$, которое
    по закону двойного отрицания есть исходное утверждение $A$.
\end{theorem}
